\begin{textsample}{POS Dim 4 – 2009\_c6 – Score -3.00 – 2009\_c6/enade\_2009\_psico\_95367afe99.txt}  \label{ex:f4_pos_327}
Capítulo 3 Análise Técnica da Prova Este capítulo tem por objetivo apresentar o desempenho dos estudantes de PSICOLOGIA no ENADE/2009.
Para isso, foram calculadas as estatísticas básicas da prova como um todo, bem como as estatísticas das questões relacionadas à Formação Geral e ao Componente Específico.
Nas Tabelas, são apresentadas as seguintes estatísticas : total da população, de presentes, média, erro-padrão da média, desvio-padrão, mediana, nota máxima, nota mínima e coeficiente de assimetria ( ver definição em 1.4.3 ).
Nas estatísticas das questões discursivas, foram detalhadas também as notas nulas : em branco, por protesto ou por insuficiência.
As estatísticas apresentadas neste capítulo contemplam, separadamente, os ingressantes, os concluintes e o total de estudantes, e foram calculadas tendo -se em vista as seguintes agregações : ( a ) Região Geográfica e País como um todo e ( b ) a Categoria Administrativa e a Organização Acadêmica.
As notas brutas das questões discursiva em cada uma das partes, Formação Geral/Componente Específico, foram calculadas a partir da média simples das notas de cada uma das questões correspondentes.
A primeira questão discursiva de Formação Geral foi desconsidera e, portanto, a nota bruta referente às questões discursiva de Formação Geral foi igual à nota da questão mantida, a segunda.
Na Componente Específica, a nota bruta correspondeu à média das notas nas três questões.
As notas brutas das questões objetivas de cada uma das partes, Formação Geral/Componente Específico, foram calculadas a partir dos acertos dos estudantes considerando -se uma regra de três para a conversão para a escala de 0 a 100, descontando -se aquelas questões que foram desconsideradas pela comissão ou que não passaram pelo teste do índice de discriminação ponto-bisserial ( ver seção 1.4.6 para a descrição do índice ).
As questões com índices de discriminação inferior a 0,2 foram excluídas do cômputo da nota do estudante.
As notas brutas de Formação Geral e do Componente Específico para cada um dos estudantes foram calculadas como médias ponderadas das suas notas nas questões objetivas e discursivas, ponderação diferente para cada parte : C\_FG = C\_FG\textasciicircum{}OBJ × 0,6 + C\_FG\textasciicircum{}DIS × 0,4 C\_CE = C\_CE\textasciicircum{}OBJ × 0,85 + C\_CE\textasciicircum{}DIS × 0,15, em que C\_FG é a nota bruta da Formação Geral, C\_FG\_OBJ é a nota bruta nas questões objetivas da Formação Geral, C\_FG\_DIS é a nota bruta nas questões discursivas da Formação Geral, C\_CE é a nota bruta do Componente Específico, C\_CE\_OBJ é a nota bruta nas questões objetivas do Componente Específico, e C\_CE\_DIS é a nota bruta nas questões discursivas do Componente Específico.
A nota bruta da prova, C\_GER, por sua vez, é também uma média ponderada da nota de cada uma das partes, Formação Geral/Componente Específico : C\_GER = C\_FG × 0,25 + C\_CE × 0,75. 3.1 Estatísticas Básicas da Prova 3.1.1 Estatísticas Básicas Gerais A Tabela 3.1 apresenta as estatísticas básicas da prova por grupo de estudantes.
O total de estudantes inscritos para a prova foi de 50.276.
Destes, 82,4 \% compareceram, sendo que a presença foi maior entre os concluintes ( 92,9 \% ), do que entre ingressantes ( 75,8 \% ).
A média geral da prova foi de 40,58, sendo que os ingressantes obtiveram média mais baixa, 35,68 que os concluintes, cuja média foi 46,98.
O desvio-padrão geral foi de 16,6, sendo o do grupo de ingressantes menor ( 14,7 ), do que o do grupo dos concluintes ( 16,8 ), aqueles possuíram uma distribuição mais homogênea nas notas.
A nota máxima foi 94,3, obtida por pelo menos um concluinte ao passo que a maior nota obtida por um ingressante foi 88,5.
Os coeficientes de assimetria para os concluintes e a população como um todo são negativos, indicando uma maior concentração à direita da distribuição e um maior espalhamento à esquerda.
Enquanto, o coeficiente de assimetria para os ingressantes foi positivo na distribuição de notas, indicando a situação oposta.
Uma comparação entre o desempenho de ingressantes e concluintes pode ser realizada a partir dos dados apresentados no Gráfico 3.1.
Cumpre destacar que as notas dos concluintes, de um modo geral, foram mais elevadas do que as dos ingressantes, cuja distribuição de notas está à esquerda dos primeiros.
A ideia de melhor desempenho é reforçada pelas notas medianas dos dois grupos : 47,9 entre os concluintes e 35,0 entre os ingressantes.
Nas faixas de notas até 50 pontos o percentual de ingressantes é maior que o dos concluintes, situação esta que se inverte para as faixas de notas superiores a 50 pontos.
Constatou -se que aproximadamente 17,4 \% dos ingressantes e 45,1 \% dos concluintes situam -se nas faixas de notas superiores a 50,0 pontos.
O Gráfico 3.2 apresenta informações referentes ao desempenho geral de ingressantes e concluintes, comparando os resultados considerando a Região, Categoria Administrativa e Organização Acadêmica.
Levando -se em conta as notas médias dos estudantes em cada Região, observa -se que, nos dois grupos, há uma variação expressiva entre as Regiões.
As notas médias mais elevadas entre os concluintes ocorreram nas Regiões Sul ( 50,08 ), Nordeste ( 47,65 ) e Centro-Oeste ( 47,53 ), todas acima da média nacional ( 46,98 ).
Em relação aos ingressantes, as médias mais altas foram obtidas nas Regiões Sul ( 37,22 ) e Nordeste ( 36,87 ), nas demais Regiões a média foi inferior à média nacional ( 35,68 ).
Na Região Norte foi obtida a média de notas mais baixa entre os ingressantes ( 31,30 ) e entre os concluintes ( 43,77 ).
Cursos nesta área foram oferecidos em instituições federais, estaduais, municipais e privadas.
Levando -se em conta os agrupamentos dos estudantes em categorias administrativas, observa -se que a pontuação média mais elevada entre os concluintes e ingressantes foi encontrada em instituições de origem federal ( respectivamente 52,43 e 43,23 ) – ambas acima da média nacional.
Tendo como foco as notas médias dos estudantes agrupados por Organização Acadêmica, notamos que universidades apresentam as maiores médias tanto entre concluintes como entre ingressantes ( respectivamente 47,61 e 36,69 ). 3.1.2 Estatísticas Básicas de Formação Geral A Tabela 3.2 apresenta as estatísticas básicas de Formação Geral por grupo de estudantes.
O total de estudantes inscritos para a prova foi de 50.276.
Destes, 82,4 \% do total compareceram, sendo que a presença foi maior entre concluintes ( 92,9 \% ) do que entre ingressantes ( 75,8 \% ).
A média geral da parte da prova de Formação Geral foi de 45,17, sendo que os ingressantes obtiveram média mais baixa, 42,98, que os concluintes, cuja média foi 48,04.
O desvio-padrão geral foi de 22,8, sendo o do grupo de ingressantes ( 22,3 ), e o do grupo dos concluintes ( 23,0 ), aqueles possuíram uma distribuição mais homogênea nas notas.
A nota máxima foi 100, obtida por estudantes de ambos os grupos.
Para esta parte, assim como para a nota geral, os coeficientes de assimetria são positivos para ingressantes e o total da população.
O coeficiente é nulo para a distribuição de notas dos concluintes.
O Gráfico 3.3 permite a comparação entre concluintes e ingressantes relativamente ao desempenho na parte da prova que aborda a Formação Geral.
Nesta parte, tal como na prova como um todo, também os concluintes apresentaram desempenho superior, embora seja perceptível um maior equilíbrio.
A mediana dos concluintes foi superior à dos ingressantes, o mesmo acontecendo com a média.
Dos ingressantes, 37,8 \% alcançaram nota superior a 50 ; dentre os concluintes, o percentual dos estudantes nesta faixa de notas foi maior, 47,7 \%.
No Gráfico 3.4 são apresentadas as informações referentes ao desempenho dos ingressantes e dos concluintes na parte de Formação Geral, em diferentes agregações como Região, Categoria Administrativa e Organização Acadêmica.
Considerando -se as notas médias dos estudantes segundo cada Região, observou -se que, tanto no grupo dos ingressantes quanto no grupo dos concluintes, houve diferenças marcantes.
As médias variam, entre os concluintes, de 42,16 na Região Norte a 52,11 na Região Sul.
Entre os ingressantes, os valores foram de 35,94 na Região Norte a 45,50 na Região Sul.
Em todas as Regiões, os concluintes apresentaram médias maiores do que os ingressantes.
Tendo como foco as categorias administrativas, as maiores notas médias foram das instituições federais, tanto para concluintes ( 53,26 ) como para ingressantes ( 50,70 ), enquanto as das instituições estaduais e municipais foram menores para os concluintes e ingressantes, respectivamente.
Os estudantes das instituições privadas se apresentam mais homogêneos do que os das demais instituições.
Considerando -se o tipo de Organização Acadêmica, Centros Universitários apresentaram os maiores valores para os concluintes ( 48,99 ) e Universidades para os ingressantes ( 44,48 ).
Os Centros Federais de Educação Tecnológica e os Institutos de Educação, Ciência e Tecnologia não estavam representados nesta área. 3.1.3 Estatísticas Básicas do Componente Específico A Tabela 3.3 apresenta as estatísticas básicas referentes ao Componente Específico da prova, tendo sua média geral, de 39,09, inferior ao desempenho da prova de PSICOLOGIA como um todo, cuja média foi 40,58, como apresentado na Tabela 3.1, que por sua vez foi inferior ao desempenho na Formação Geral, que teve média igual a 45,17, conforme mostrado na Tabela 3.2.
Os concluintes apresentaram um desempenho médio ( 46,68 ) superior ao dos ingressantes ( 33,30 ), semelhante ao ocorrido para a nota da Formação Geral, havendo ainda indícios de maior homogeneidade entre as notas médias dos ingressantes do que entre as dos concluintes, como sugeriu o desvio padrão, menor no grupo dos ingressantes do que no dos concluintes.
As notas máximas obtidas por concluintes e ingressantes foram iguais a 93,1 e 94,9, respectivamente.
As distribuições de notas para o grupo dos concluintes se apresentaram com uma concentração nas notas maiores com cauda mais pesada nas notas menores, como evidenciados pelo coeficiente de assimetria negativo.
A situação inversa ocorre para os ingressantes.
Estes coeficientes, porém, são pequenos.
O total da população apresentou um coeficiente nulo, caracterizando uma simetria da curva de distribuição das notas.
A situação de assimetria é confirmada pelo Gráfico 3.5, que proporciona uma comparação do desempenho entre concluintes e ingressantes em relação ao Componente Específico da prova.
Claramente a distribuição das notas médias dos ingressantes, à direita daquela dos concluintes, evidenciou o percentual mais alto de notas médias nas faixas mais baixas para este grupo.
As medianas das distribuições das notas dos ingressantes e dos concluintes foram 32,3 e 48,0, respectivamente, enfatizando a superioridade dos concluintes nesta parte da prova.
Assim como os Gráficos 3.2 e 3.4, o Gráfico 3.6, apresenta uma comparação dos resultados em relação às Regiões, Categoria Administrativa e Organização Acadêmica, agora levando em conta o desempenho dos ingressantes e dos concluintes no Componente Específico da prova.
Tendo como foco as notas médias dos estudantes, em determinada Região, observa -se que, entre os concluintes, as mais elevadas foram encontradas nas Regiões Sul ( 49,45 ), Centro Oeste ( 48,30 ) e Nordeste ( 47,89 ).
Entre os ingressantes, as maiores notas ocorreram também nas mesmas Regiões : Nordeste ( 35,24 ), Sul ( 34,52 ) e Centro-Oeste ( 33,57 ).
Novamente, para a análise segundo a Categoria Administrativa, todas as esferas foram representadas.
As médias mais elevadas, tanto entre os concluintes como entre os ingressantes, foram encontradas em instituições de origem federal ( 52,20 entre os concluintes e 40,79 entre os ingressantes ) enquanto os valores mais baixos foram encontrados nas instituições Estaduais ( concluintes 37,41 ) e Municipais ( ingressantes 27,11 ).
Quanto à Organização Acadêmica, as Universidades, os Centros Universitários e as Faculdades foram representados nesta área.
Considerando -se o tipo de Organização Acadêmica, a ordenação das médias foram as mesmas para concluintes e ingressantes : Universidades ( 47,24 e 34,14 respectivamente ) seguidas dos Centros Universitários ( 45,79 e 32,84 respectivamente ) e das Faculdades ( 45,21 e 31,97 respectivamente ) 3.2 Análise das Questões Objetivas 3.2.1 Formação Geral A Tabela 3.4 apresenta as estatísticas básicas relativas às oito questões objetivas da Formação Geral dos estudantes.
Adiantamos que a primeira questão foi desconsiderada pelo critério da correlação ponto-bisserial.
Como pode ser observado na Tabela 3.4, os estudantes concluintes obtiveram desempenho superior, com nota média de 57,94, em relação aos estudantes ingressantes cuja nota média foi 53,12.
A variabilidade entre os concluintes foi ligeiramente maior do que entre os ingressantes, como indicaram os desvios-padrão de 24,1, para o grupo de ingressantes e 24,4 para o grupo dos concluintes.
A amplitude total das notas foi a mesma para ambos os grupos : nota mínima de 0 e máxima de 100.
A mediana foi igual para ambos os grupos : 57,1.
A assimetria, medida pelo coeficiente de assimetria ( skewness ) foi um pouco menor em módulo para os ingressantes ( -0,2 ) do que para os concluintes ( -0,4 ), ainda que ambos fossem negativos indicando uma assimetria com cauda para a esquerda e concentração à direita.
A Tabela 3.5 apresenta a classificação das questões objetivas da Formação Geral, segundo o índice de facilidade ( questões em vermelho foram desconsideradas pelo critério ponto-bisserial ).
Como as questões foram comuns a todas as áreas, a tabela apresenta a classificação considerando -se os concluintes de todas as áreas e aqueles do curso de PSICOLOGIA.
De acordo com os índices obtidos, as questões objetivas da prova foram assim avaliadas para a área de PSICOLOGIA : das oito questões, três tiveram entre 60 \%, inclusive, e 85 \% de acertos e foram classificadas como fáceis ( somente uma para o total de concluintes ) ; outras duas foram consideradas medianas, alcançando entre 40 \%, inclusive, e 60 \% de acertos ( três para o total de concluintes ) ; duas obtiveram entre 15 \%, inclusive, e 40 \% de acertos, sendo consideradas difíceis ( três para o total de concluintes ).
Por fim, uma foi classificada como muito difícil tento obtido percentual menor do que 15 \% de acertos ( também uma para o total de concluintes ), coincidentemente a questão desconsiderada pelo critério ponto-bisserial.
Nenhuma questão foi classificada na categoria Muito fácil.
Em linhas gerais, podemos dizer que os concluintes de PSICOLOGIA se deram melhor nesta parte em relação ao total da população de concluintes de todos os cursos.
A Tabela 3.6, descreve os resultados das análises das questões objetivas relativas à Formação Geral, segundo o poder de discriminação, utilizando -se para tal o índice de discriminação ponto-bisserial.
Como as questões foram comuns a todas as áreas, a tabela apresenta a classificação considerando -se os concluintes de todas as áreas e aqueles do curso de PSICOLOGIA.
A desconsideração da questão segundo este critério considerou o universo de todos os concluintes, três das oito questões apresentaram índices igual ou acima de 0,40 ( nenhuma para a população de concluintes como um todo ), sendo, portanto, classificadas como muito boas nesta característica, para este grupo de estudantes, nenhuma foi classificada como boa com índice entre 0,30, inclusive, e 0,40 ( três para os concluintes como um todo ) e quatro foram classificadas como médias, com índice entre 0,20, inclusive, e 0,30 ( também quatro para os concluintes como um todo ).
Uma das questões foi classificada como fraca tanto para a população de concluintes como um todo como para a de concluintes de PSICOLOGIA e foi desconsiderada pelo critério bisserial. concluintes ) ; outras duas foram consideradas medianas, alcançando entre 40 \%, inclusive, e 60 \% de acertos ( três para o total de concluintes ) ; duas obtiveram entre 15 \%, inclusive, e 40 \% de acertos, sendo consideradas difíceis ( três para o total de concluintes ).
Por fim, uma foi classificada como muito difícil tento obtido percentual menor do que 15 \% de acertos ( também uma para o total de concluintes ), coincidentemente a questão desconsiderada pelo critério ponto-bisserial.
Nenhuma questão foi classificada na categoria Muito fácil.
Em linhas gerais, podemos dizer que os concluintes de PSICOLOGIA se deram melhor nesta parte em relação ao total da população de concluintes de todos os cursos.
A Tabela 3.6, descreve os resultados das análises das questões objetivas relativas à Formação Geral, segundo o poder de discriminação, utilizando -se para tal o índice de discriminação ponto-bisserial.
Como as questões foram comuns a todas as áreas, a tabela apresenta a classificação considerando -se os concluintes de todas as áreas e aqueles do curso de PSICOLOGIA.
A desconsideração da questão segundo este critério considerou o universo de todos os concluintes, três das oito questões apresentaram índices igual ou acima de 0,40 ( nenhuma para a população de concluintes como um todo ), sendo, portanto, classificadas como muito boas nesta característica, para este grupo de estudantes, nenhuma foi classificada como boa com índice entre 0,30, inclusive, e 0,40 ( três para os concluintes como um todo ) e quatro foram classificadas como médias, com índice entre 0,20, inclusive, e 0,30 ( também quatro para os concluintes como um todo ).
Uma das questões foi classificada como fraca tanto para a população de concluintes como um todo como para a de concluintes de PSICOLOGIA e foi desconsiderada pelo critério bisserial.
Na sequência, no Quadro 3.1, estão especificadas, por questão, as habilidades aferidas pelas questões de Formação Geral da prova da PSICOLOGIA, bem como os resultados da análise relativa aos índices de facilidade e de discriminação de cada questão.
O índice de facilidade para a área de PSICOLOGIA teve uma variação considerável, oscilando entre 0,10 até 0,82 ( na população de concluintes o intervalo foi de 0,08 a 0,75 ) e o de discriminação ficou entre 0,07 e 0,47 ( na população de concluintes o intervalo foi de 0,02 a 0,39 ).
A questão 8, que aferiu as habilidades de “ Ler e interpretar textos ” ; “ Analisar e criticar informações ” ; “ Extrair conclusões por indução e/ou dedução ” ; “ Estabelecer relações, comparações e contrastes em diferentes situações ” ; “ Detectar contradições ” ; “ Fazer escolhas valorativas avaliando consequências ” ; “ Questionar a realidade ” ; “ Argumentar coerentemente ” foram as que apresentaram, para a população de concluintes de PSICOLOGIA, maior poder discriminatório, com índice de 0,47.
A questão mais fácil, com 82 \% de acertos, foi também a de número 8.
De outra parte, a questão com menor índice de facilidade e de discriminação em Formação Geral foi a questão 1, que foi desconsiderada, e que deveria aferir as habilidades de “ Ler e interpretar textos ” ; “ Analisar e criticar informações ” ; “ Extrair conclusões por indução e/ou dedução ” ; “ Estabelecer relações, comparações e contrastes em diferentes situações ” ; “ Detectar contradições ” ; “ Fazer escolhas valorativas avaliando consequências ” ; “ Questionar a realidade ” ; “ Argumentar coerentemente ”.
O Gráfico 3.7 mostra, a título de ilustração, o comportamento das respostas à questão 2 de Formação Geral os estudantes concluintes da área de PSICOLOGIA como função da nota geral na prova.
Trata -se de uma questão considerada fácil para os concluintes de PSICOLOGIA e média para o conjunto de concluintes de todas as áreas.
O poder de discriminação foi médio para os estudantes da área e médio para o conjunto de todos os concluintes.
Neste Gráfico, cada uma das cinco curvas, representou o percentual de respostas em determinada alternativa da questão, em função do escore total dos estudantes concluintes da área de PSICOLOGIA.
A curva vermelha corresponde à alternativa C, a correta para esta questão.
Assim, observa -se que entre os estudantes com notas mais baixas – desempenho geral mais fraco na prova como um todo –, a tendência foi, além da escolha da alternativa correta, a escolha de alternativas incorretas, especialmente a alternativa E – curva roxa, seguida da alternativa D – curva em verde escuro.
A proporção de estudantes que selecionaram a resposta correta C, aumentou com o desempenho deste grupo de estudantes chegando ao acerto de 100 \% dos estudantes com notas pelo menos 87,5.
Esta análise permitiu verificar como a questão discriminou entre os grupos de desempenho.
Os Gráficos relativos às demais questões de Formação Geral constam do Anexo I. 3.2.2 Componente Específico A Tabela 3.7 apresenta as estatísticas básicas em relação às questões objetivas do Componente Específico.
Como pode ser observado, os estudantes concluintes obtiveram desempenho superior, com uma média de 50,04 pontos, em relação aos estudantes ingressantes com uma média de 36,89 pontos.
Os desvios-padrão indicaram que a variabilidade entre os concluintes foi superior ( 18,5 ) do que entre os ingressantes ( 16,4 ).
O coeficiente de assimetria para os ingressantes foi positivo indicando uma leve inclinação para à esquerda enquanto que para os concluintes foi negativo, indicando uma inclinação para a direita.
A assimetria foi, porém, maior em módulo entre os concluintes.
A análise das questões objetivas do Componente Específico em PSICOLOGIA, em relação ao grau de facilidade da prova ( ver Tabela 3.8 ), mostrou que, das 27 questões, 14 encontraram -se na categoria difícil ( entre 15 \% e 40 \%, não incluso, de acertos ), 10 foram de nível médio ( entre 40 \% e 60 \%, não incluso, de acertos ), 2 estão na categoria fácil ( entre 60 \% e 85 \%, não incluso, de acertos ) e nenhuma teve a classificação muito fácil ( a partir de 85 \% de acertos ).
Nenhuma questão teve percentual de acertos menor do que 15 \% ( classificação muito difícil ) e 1 questão foi desconsiderada pela Comissão Assessora da Área de PSICOLOGIA.
Das 27 questões objetivas relativas ao Componente Específico ( ver Tabela 3.9 ), 7 obtiveram o índice de discriminação muito bom, 7 obtiveram o índice de discriminação bom e 6 obtiveram nível de discriminação médio.
O índice fraco que eliminaria a questão, foi alcançado por 6 questões e uma questão foi desconsiderada pela comissão.
Na sequência, no Quadro 3.2 estão especificadas, por questão, as habilidades aferidas pelas questões do Componente Específico da prova, bem como os resultados da análise relativa aos índices de facilidade e de discriminação de cada questão.
Das 27 questões objetivas relativas ao Componente Específico ( ver Tabela 3.9 ), 7 obtiveram o índice de discriminação muito bom, 7 obtiveram o índice de discriminação bom e 6 obtiveram nível de discriminação médio.
O índice fraco que eliminaria a questão, foi alcançado por 6 questões e uma questão foi desconsiderada pela comissão.
Na sequência, no Quadro 3.2 estão especificadas, por questão, as habilidades aferidas pelas questões do Componente Específico da prova, bem como os resultados da análise relativa aos índices de facilidade e de discriminação de cada questão.
O Gráfico 3.8 mostra, a título de ilustração, o comportamento da questão 12 do Componente Específico para os estudantes concluintes da área de PSICOLOGIA como função da nota geral na prova.
Trata -se de uma questão considerada média para os concluintes de PSICOLOGIA.
O poder de discriminação foi bom para os estudantes da área.
Neste Gráfico, cada uma das cinco curvas representa o percentual de respostas em determinada alternativa da questão, em função da nota total dos estudantes concluintes de PSICOLOGIA.
A curva vermelha correspondeu à alternativa C, a correta para esta questão.
Assim, observou -se que entre os estudantes com escores mais baixos – desempenho geral mais fraco na prova como um todo –, ocorreram escolhas de uma alternativa incorreta como a alternativa B – curva verde claro ou a alternativa E – curva roxa.
A proporção de estudantes que selecionaram a resposta correta aumentou com o desempenho deste grupo de estudantes chegando ao acerto de 100 \% dos estudantes com notas pelo menos 87,5.
Esta análise permitiu verificar como a questão discriminou entre os grupos de desempenho.

% matched lemmas: 
\end{textsample}
